\documentclass[leqno]{article}
% Shared mark scheme preamble for resources.danieljsmith.org: the pack
% preamble, unchanged, then the scheme's own layer. A scheme's .tex is
% rendered from its YAML (scripts/README.md, Mark schemes), never edited.
\input{\CurrentFilePath/pack_preamble.tex}

% The page: equal side margins, so the table runs the full width of the
% header rule. The footers' right edge follows the text's.
\geometry{left=0.8in, right=0.8in}
\makeatletter
\renewcommand{\djs@footR}[1]{%
  \put(\dimexpr1in+\hoffset+\oddsidemargin+\textwidth\relax,
    -\dimexpr\paperheight-0.32in\relax){\makebox[0pt][r]{#1}}}
\makeatother

\usepackage[export]{adjustbox}
\usepackage{tabularray}
\definecolor{msGrid}{RGB}{176, 172, 164}
\definecolor{msTint}{RGB}{245, 243, 239}

% The header on every page, the front page included, and the booklet label
% under the contents-page title. \djsQuestionPageStyle stays empty.
\newcommand{\djsSchemeHeader}[1]{%
  \djsPackHeader[markscheme]{\textit{Mark scheme} - #1}%
  \renewcommand{\djsPackTitle}{#1}%
  \renewcommand{\djsBookletLabel}{Mark scheme}%
}

% The front page: the contents heading, spaced as tocloft spaces the packs'
% "Questions", then the contents table. A question's name and page number
% both link to its first page.
\newcommand{\msContentsHeading}{%
  \vskip\cftbeforetoctitleskip
  {\noindent\cfttoctitlefont\contentsname\par}%
  \vskip\cftaftertoctitleskip}
\NewTblrEnviron{mscontents}
\SetTblrInner[mscontents]{
  colspec = {Q[wd=26mm, l] X[l] Q[wd=16mm, c] Q[wd=14mm, r]},
  width = \linewidth,
  rowsep = 4pt, colsep = 6pt,
  hlines = {0.4pt, msGrid}, vlines = {0.4pt, msGrid},
  hline{2} = {2.2pt, scarlet},
  row{1} = {font=\small\bfseries},
  row{Z} = {font=\bfseries},
  column{2} = {font=\small},
}
\newcommand{\msContentsQuestion}[1]{\hyperref[q:#1]{Question #1}}
\newcommand{\msContentsPage}[1]{\hyperref[q:#1]{\pageref*{q:#1}}}
\newcommand{\msSep}{\hspace{0.9em}}

% The conventions pages' layout: headings in the documents' own bold, tables
% ruled and padded as the scheme's, roomy lists.
\newcommand{\msHead}[1]{%
  \par\vspace{18pt}%
  {\noindent\bfseries #1\par}%
  \vspace{6pt}}
\NewTblrEnviron{mstblr}
\SetTblrInner[mstblr]{rowsep = 6pt, colsep = 6pt,
  hlines = {0.4pt, msGrid}, vlines = {0.4pt, msGrid}}
\newenvironment{msList}
  {\begin{itemize}[leftmargin=1.2em, itemsep=5pt, topsep=3pt, parsep=0pt]}
  {\end{itemize}}

% A question's first page: the link target and the PDF bookmark, as
% \questionitem gives the packs.
\newcommand{\msQnum}{}
\newcommand{\msQuestion}[1]{%
  \clearpage
  \renewcommand{\msQnum}{#1}%
  \phantomsection\label{q:#1}%
  \addcontentsline{toc}{subsection}{\protect\djsTocEntry{#1}}}

% The scheme's table: Question 16mm, Scheme the rest, Marks 16mm, Notes
% 44mm; every cell ruled in msGrid; the scarlet rule under the head row. A
% long table has no caption or "continued" foot; a continuation page opens
% with a muted line above the repeated head row. The rules across a part's
% rows, the alternative's tint and the total row are the rendered table's.
\NewTblrTheme{ms}{%
  \DefTblrTemplate{firsthead}{default}{}%
  \DefTblrTemplate{middlehead,lasthead}{default}{%
    {\small\itshape\color{djsMuted}Question \msQnum\ continued\par}%
    \vspace{6pt}}%
  \DefTblrTemplate{firstfoot,middlefoot,lastfoot}{default}{}%
}
\NewTblrEnviron{msscheme}
\SetTblrOuter[msscheme]{long, theme=ms, entry=none, label=none}
\SetTblrInner[msscheme]{
  colspec = {Q[wd=16mm, l] X[l] Q[wd=16mm, c] Q[wd=44mm, l]},
  width = \linewidth,
  rowsep = 6pt, colsep = 6pt,
  rowhead = 1,
  vlines = {0.4pt, msGrid},
  hline{1} = {0.4pt, msGrid}, hline{2} = {2.2pt, scarlet},
  hline{Z} = {0.4pt, msGrid},
  column{1} = {font=\bfseries}, column{3} = {font=\bfseries},
  column{4} = {font=\small},
  row{1} = {rowsep=4pt, font=\small\bfseries},
}

% Two marks on a row: each note starts with its code.
\newcommand{\msNote}[2]{\textbf{#1}\enspace #2}
% A full-row figure: the worked solution's own picture, scaled uniformly and
% only down, to the table less the Question column. The picture is drawn
% with the line width it has in the solutions document, at its list depth
% there (the optional argument, 1 by default: a part's working), so one
% sized by \linewidth keeps its size; tabularray sets the cell's \linewidth
% to the cell's measured width, which would shrink it at every pass. The
% widths are measured once, in a solution box at the packs' text width
% (pack_preamble.tex's margins, 0.8in and 1in).
\newlength{\msSpanW}
\makeatletter
\newcommand{\ms@keepLW}[1]{%
  \expandafter\xdef\csname ms@solnLW@#1\endcsname{\the\linewidth}}
\AtBeginDocument{%
  \setlength{\msSpanW}{\dimexpr\textwidth-16mm-28pt\relax}%
  \setbox\z@\vbox{%
    \begin{minipage}{\dimexpr\paperwidth-1.8in\relax}
      \begin{djsSolution}
        \ms@keepLW{0}%
        \begin{enumerate}\item \ms@keepLW{1}%
        \begin{enumerate}\item \ms@keepLW{2}%
        \begin{enumerate}\item \ms@keepLW{3}%
        \end{enumerate}\end{enumerate}\end{enumerate}
      \end{djsSolution}
    \end{minipage}}}
\newcommand{\msFigure}[2][1]{%
  \adjustbox{max width=\msSpanW, max totalheight=0.6\textheight, valign=t}{%
    \hsize=\csname ms@solnLW@#1\endcsname
    \linewidth=\hsize
    #2}}
\makeatother
% A figure mark's feature, in the Scheme column.
\newcommand{\msFeature}[1]{\textit{#1}}
% An alternative's title row.
\newcommand{\msAlternative}[1]{\textit{Alternative:} #1}
% Drafts only: a queried mark's code, the query's why under its note, and a
% flag in its part's place.
\newcommand{\msQueried}[2]{#1~\(\to\) #2?}
\newcommand{\msQuery}[1]{{\itshape\color{djsMuted}Query: #1}}
\newcommand{\msFlag}[2]{\textbf{#1:}\enspace #2}


\djsSchemeHeader{Volumes of Revolution}
\djsSitePath{Core Pure}{Further Calculus}

\djsLastUpdated{25 October 2026}
\begin{document}
\djsSetQuestionCount{15}
\djsAddMarks{1}{6}
\djsAddMarks{2}{7}
\djsAddMarks{3}{5}
\djsAddMarks{4}{5}
\djsAddMarks{5}{10}
\djsAddMarks{6}{8}
\djsAddMarks{7}{5}
\djsAddMarks{8}{7}
\djsAddMarks{9}{6}
\djsAddMarks{10}{8}
\djsAddMarks{11}{6}
\djsAddMarks{12}{7}
\djsAddMarks{13}{9}
\djsAddMarks{14}{11}
\djsAddMarks{15}{9}
\vspace*{6pt}
\djsContentsTitle
\vspace{30pt}
\msContentsHeading
\begin{mscontents}{}
Question & Marks by part & Marks & Page \\
\msContentsQuestion{1} & (a)\,2\msSep (b)\,4 & 6 & \msContentsPage{1} \\
\msContentsQuestion{2} & (a)\,2\msSep (b)\,5 & 7 & \msContentsPage{2} \\
\msContentsQuestion{3} & & 5 & \msContentsPage{3} \\
\msContentsQuestion{4} & & 5 & \msContentsPage{4} \\
\msContentsQuestion{5} & (a)\,2\msSep (b)\,5\msSep (c)\,3 & 10 & \msContentsPage{5} \\
\msContentsQuestion{6} & & 8 & \msContentsPage{6} \\
\msContentsQuestion{7} & & 5 & \msContentsPage{7} \\
\msContentsQuestion{8} & & 7 & \msContentsPage{8} \\
\msContentsQuestion{9} & & 6 & \msContentsPage{9} \\
\msContentsQuestion{10} & (a)\,6\msSep (b)\,2 & 8 & \msContentsPage{10} \\
\msContentsQuestion{11} & & 6 & \msContentsPage{11} \\
\msContentsQuestion{12} & & 7 & \msContentsPage{12} \\
\msContentsQuestion{13} & (a)\,3\msSep (b)\,6 & 9 & \msContentsPage{13} \\
\msContentsQuestion{14} & (a)\,4\msSep (b)\,7 & 11 & \msContentsPage{14} \\
\msContentsQuestion{15} & (a)\,3\msSep (b)\,6 & 9 & \msContentsPage{15} \\
\SetCell[c=2]{r} Total & & 109 & \\
\end{mscontents}
\newpage
% The conventions pages, shared by every pack's mark scheme: one source of
% truth, \input after the front page, with no strand or pack logic. The keys
% on the first side, the rules on the second. The scheme layer defines
% \msHead, msList and mstblr.
\vspace*{6pt}%
{\noindent\cfttoctitlefont Reading the mark scheme\par}%
\vspace{10pt}%
These rules apply throughout, unless a note says otherwise.

\msHead{Marks}
Every mark is worth one, so \textbf{M1 A1} on a row is two marks. Only marks
coded ft follow through.\par
\vspace{10pt}%
\begin{mstblr}{
  colspec = {Q[wd=15mm, l] X[l]},
  width = \linewidth,
  hline{2} = {2.2pt, scarlet},
  column{1} = {font=\bfseries},
  row{1} = {font=\small\bfseries, rowsep=4pt},
}
Code & Meaning \\
M1 & Method: a correct method, applied to this question. A slip doesn't cost it if the method is clear. A formula must be used, not just quoted. \\
dM1 & Dependent method: as M1, but only if the method mark before it was earned, unless its note names another. \\
A1 & Accuracy: a correct answer or step. It needs the method mark it follows. \\
A1ft & As A1, following through from an earlier answer that the note names. \\
A1~AG & The last mark of a part whose answer is printed in the question: complete working, with no errors. \\
B1 & A correct value, fact or statement. It needs no method mark. \\
B1ft & As B1, following through from an earlier answer that the note names. \\
E1 & A written explanation, interpretation or conclusion. Concluding a proof by induction needs every earlier mark. \\
\end{mstblr}

\msHead{Abbreviations}
ft and AG qualify a code in the Marks column. The others appear in the
notes.\par
\vspace{10pt}%
\begin{mstblr}{
  colspec = {Q[wd=17mm, l] X[l] Q[wd=15mm, l] X[l]},
  width = \linewidth,
  column{1} = {font=\bfseries}, column{3} = {font=\bfseries},
}
ft & follow through & AG & answer given in the question \\
cao & correct answer only, with no follow-through & cso & correct solution only: no errors in the part \\
oe & or an equivalent form & awrt & any answer that rounds to this value \\
soi & seen, or implied by later working & isw & ignore working after a correct answer \\
SC & special case: credit outside the main method & their & the student's own earlier value \\
condone & a slip that doesn't cost the mark & [\;\,] & working the mark doesn't need \\
s.f. & significant figures & d.p. & decimal places \\
\end{mstblr}

\newpage
\msHead{Reading the scheme}
\begin{msList}
\item Each row is a line of working, with the marks it earns beside it.
  Where a row carries two marks, each note starts with its code.
\item \textbf{[\textit{n}]} ends each part with its total.
\item A shaded \textit{Alternative} is a second method for the same total.
  Mark an answer by the method it follows.
\item Any other valid method earns the equivalent marks, unless the question
  names a method or says \textit{hence}.
\item A sketch earns one B1 for each feature named below its figure.
\end{msList}

\msHead{Accuracy and answers}
\begin{msList}
\item Unless a question says otherwise, give answers to 3 significant
  figures, and angles in degrees to 1 decimal place. Angles in radians take
  3 significant figures.
\item A final answer with more figures is accepted if it rounds to the
  scheme's value; one with fewer is not. An answer spoiled by rounding too
  early loses its accuracy mark.
\item Where a question asks for an exact answer, a decimal loses the final
  accuracy mark. Otherwise either is accepted.
\item Wrong or missing units cost nothing unless the question asks for them.
\item A correct answer with no working earns full marks, unless the question
  needs the working: show that, prove, hence, a method it names, or an exact
  answer. A wrong answer with no working earns nothing. Calculator answers
  are treated the same way.
\item A correct answer, once reached, isn't spoiled by later working (isw).
\end{msList}

\msHead{Value of $\boldsymbol{g}$}
\begin{msList}
\item Take $g = 9.8\,\mathrm{m\,s^{-2}}$ unless a question states a value,
  which is then the only one accepted. Answers that use $g = 9.8$ may be
  given to 2 or 3 significant figures.
\item Where a question doesn't state $g$, an answer worked with $g = 10$
  throughout earns the same marks, keeping any value printed in the
  question. The worked solutions list the final answers with $g = 10$.
\end{msList}


\msQuestion{1}
\begin{msscheme}{
  hline{3-4,6-9} = {2-Z}{0.4pt, msGrid},
  hline{5,10} = {0.4pt, msGrid},
  row{10} = {font=\bfseries},
}
Question & Scheme & Marks & Notes \\
1(a) & {\(\frac{\mathrm{d}}{\mathrm{d}x}\ln(\sec x+\tan x) = \frac{\sec x\tan x+\sec^2 x}{\sec x+\tan x}\)} & M1 & {Divides a derivative of \(\sec x+\tan x\) of the form \(p\sec x\tan x+q\sec^2 x\) by \(\sec x+\tan x\), oe in sines and cosines.} \\
& {\(\phantom{\frac{\mathrm{d}}{\mathrm{d}x}} = \frac{\sec x(\tan x+\sec x)}{\sec x+\tan x} = \sec x\)} & A1~AG & {Correct derivative simplified without error to \(\sec x\); combined steps accepted.} \\
\nopagebreak & & [2] & \\
1(b) & {\(V = \pi\int_0^{\pi/3}\sec x\,\mathrm{d}x\)} & B1 & {Needs the \(\pi\); limits may be implied by their substitution.} \\
& {\(\phantom{V} = \pi\left[\ln(\sec x+\tan x)\right]_0^{\pi/3}\)} & M1 & {Obtains \(\ln(\sec x+\tan x)\), or another correct integral of \(\sec x\). The \(\pi\) need not be seen.} \\
& {\(\phantom{V} = \pi\left(\ln(2+\sqrt{3}) - \ln 1\right)\)} & dM1 & {Substitutes both limits and subtracts, evaluating \(\sec\) and \(\tan\) exactly. Condone the zero lower-limit term unwritten.} \\
& {\(V = \pi\ln(2+\sqrt{3})\)} & A1 & {oe exact, such as \(-\pi\ln(2-\sqrt{3})\).} \\
\nopagebreak & & [4] & \\
\nopagebreak \SetCell[c=2]{r} Total for Question 1 & & 6 & \\
\end{msscheme}

\msQuestion{2}
\begin{msscheme}{
  hline{3-4,6-9} = {2-Z}{0.4pt, msGrid},
  hline{5,10} = {0.4pt, msGrid},
  row{10} = {font=\bfseries},
}
Question & Scheme & Marks & Notes \\
2(a) & {\(\sinh(2\ln 2) = \frac{e^{2\ln 2} - e^{-2\ln 2}}{2}\)\\[4pt] \(e^{2\ln 2} = e^{\ln 4} = 4\)} & M1 & {Uses the exponential definition of \(\sinh\) with \(2\ln 2\), or a double-angle identity with \(\ln 2\), and evaluates at least one exponential.} \\
& {\(\sinh(2\ln 2) = \frac{4 - \frac14}{2} = \frac{15}{8}\)} & A1~AG & {Correct substituted expression leading without error to \(\frac{15}{8}\); combined steps accepted.} \\
\nopagebreak & & [2] & \\
2(b) & {\(V = \pi\int_0^{2\ln 2} (\cosh x - 1)\,\mathrm{d}x\)} & B1 & {Correct integral with limits \(0\) and \(2\ln 2\), oe such as a cylinder of volume \(2\pi\ln 2\) subtracted. Needs the \(\pi\).} \\
& {\(\int (\cosh x - 1)\,\mathrm{d}x = \sinh x - x\)} & M1 A1 & {\msNote{M1}{Integrates at least one term correctly, such as \(\cosh x \to \sinh x\).}\\ \msNote{A1}{With or without \(\pi\).}} \\
& {\(\pi\left[\sinh x - x\right]_0^{2\ln 2} = \pi\left(\tfrac{15}{8} - 2\ln 2 - 0\right)\)} & dM1 & {Applies their limits to their integral of the form \(p\sinh x + qx\), evaluating \(\sinh(2\ln 2)\) by (a) or otherwise.} \\
& {\(V = \pi\left(\frac{15}{8} - 2\ln 2\right)\)} & A1 & {Any exact equivalent, such as \(\frac{15\pi}{8} - \pi\ln 4\).} \\
\nopagebreak & & [5] & \\
\nopagebreak \SetCell[c=2]{r} Total for Question 2 & & 7 & \\
\end{msscheme}

\msQuestion{3}
\begin{msscheme}{
  hline{3-6} = {2-Z}{0.4pt, msGrid},
  hline{7} = {0.4pt, msGrid},
  row{7} = {font=\bfseries},
}
Question & Scheme & Marks & Notes \\
3 & {\(\pi\int_0^k \frac{k-x}{k+x}\,\mathrm{d}x = 1\)} & B1 & {Needs \(\pi\), \(y^2\) as the integrand and limits 0 and \(k\); the \(= 1\) may be seen later.} \\
& {\(\frac{k-x}{k+x} = -1 + \frac{2k}{k+x}\)\\[4pt] \(\Rightarrow \int \ldots \,\mathrm{d}x = \ldots\)} & M1 & {Writes the integrand as a constant plus \(\frac{c}{k+x}\) and integrates to the form \(px + q\ln(k+x)\). Condone sign and coefficient slips.} \\
& {\(\left[-x + 2k\ln(k+x)\right]_0^k\)} & A1 & {Accept \(\ln|k+x|\) or any equivalent, such as \(2k\ln(k+x) - x + c\).} \\
& {\(\pi\left(-k + 2k\ln 2k - 2k\ln k\right) = 1\)\\[4pt] \(\Rightarrow \pi k(2\ln 2 - 1) = 1 \Rightarrow k = \ldots\)} & dM1 & {Substitutes both limits into their integral, subtracts, and combines the logarithms to a constant multiple of \(k\).} \\
& {\(k = \frac{1}{\pi(2\ln 2 - 1)}\)} & A1 & {Or any exact equivalent, such as \(\frac{1}{\pi(\ln 4 - 1)}\).} \\
\nopagebreak \SetCell[c=2]{r} Total for Question 3 & & 5 & \\
\end{msscheme}

\msQuestion{4}
\begin{msscheme}{
  hline{3-6} = {2-Z}{0.4pt, msGrid},
  hline{7} = {0.4pt, msGrid},
  row{7} = {font=\bfseries},
}
Question & Scheme & Marks & Notes \\
4 & {\(\pi\int_0^k \frac{x^2}{k^2+x^2}\,\mathrm{d}x = 1\)} & B1 & {Needs the \(\pi\) and limits 0 and \(k\); the equation with 1 may appear later.} \\
& {\(\frac{x^2}{k^2+x^2} = 1 - \frac{k^2}{k^2+x^2}\)\\[4pt] \(\Rightarrow \int \ldots\,\mathrm{d}x = x - \ldots\arctan\frac{x}{k}\)} & M1 & {Writes the integrand as \(1 - \frac{p}{k^2+x^2}\) and integrates to \(x \pm q\arctan\frac{x}{k}\), or substitutes \(x = k\tan\theta\) and integrates \(\tan^2\theta\) via \(\sec^2\theta - 1\). Condone coefficient slips.} \\
& {\(\left[x - k\arctan\frac{x}{k}\right]_0^k\)} & A1 & {oe, such as \(k(\tan\theta - \theta)\) after \(x = k\tan\theta\).} \\
& {\(\pi\left(k - k\arctan 1\right) = 1 \Rightarrow k = \ldots\)} & dM1 & {Substitutes both limits, in the variable used, into their integral, equates \(\pi\) times the result to 1 and solves for \(k\).} \\
& {\(k = \frac{4}{\pi(4-\pi)}\)} & A1 & {Any exact equivalent, such as \(\frac{1}{\pi\left(1 - \frac{\pi}{4}\right)}\).} \\
\nopagebreak \SetCell[c=2]{r} Total for Question 4 & & 5 & \\
\end{msscheme}

\msQuestion{5}
\begin{msscheme}{
  hline{3-4,6-10,12-14} = {2-Z}{0.4pt, msGrid},
  hline{5,11,15} = {0.4pt, msGrid},
  row{15} = {font=\bfseries},
}
Question & Scheme & Marks & Notes \\
5(a) & {\(\int \sin^2\theta\,\mathrm{d}\theta = \int \tfrac12(1 - \cos 2\theta)\,\mathrm{d}\theta\)\\[4pt] \(\phantom{\int \sin^2\theta\,\mathrm{d}\theta} = \ldots\)} & M1 & {Uses \(\sin^2\theta = \frac12(1 \pm \cos 2\theta)\) and integrates to the form \(p\theta + q\sin 2\theta\). Condone sign and coefficient slips.} \\
& {\(= \frac{\theta}{2} - \frac{\sin 2\theta}{4} + c\)} & A1 & {oe, with \(+c\).} \\
\nopagebreak & & [2] & \\
5(b) & {\(\frac{\mathrm{d}y}{\mathrm{d}\theta} = -\operatorname{cosec}^2\theta\)} & B1 & \\
& {\(V = \pi\int (6\sin^2\theta)^2 \times \text{"}{-\operatorname{cosec}^2\theta}\text{"}\,\mathrm{d}\theta\)} & M1 & {Uses \(\pi\int x^2\frac{\mathrm{d}y}{\mathrm{d}\theta}\,\mathrm{d}\theta\) with \(x = 6\sin^2\theta\) and their \(\frac{\mathrm{d}y}{\mathrm{d}\theta}\) of the form \(\pm p\operatorname{cosec}^2\theta\). Limits not needed.} \\
& {\(= 36\pi\int -\sin^2\theta\,\mathrm{d}\theta\)} & A1 & {oe. Limits not needed.} \\
& {\(y = \sqrt3 \Rightarrow \theta = \frac{\pi}{6}\)\\[4pt] \(y = \frac{1}{\sqrt3} \Rightarrow \theta = \frac{\pi}{3}\)} & B1 & {Both, matched to their \(y\) values; may be seen as limits.} \\
& {\(V = 36\pi\int_{\pi/6}^{\pi/3} \sin^2\theta\,\mathrm{d}\theta\)\\[4pt] \(k = 36\pi,\ a = \frac{\pi}{6},\ b = \frac{\pi}{3}\)} & A1 & {cso. Accept \(k = -36\pi\), \(a = \frac{\pi}{3}\), \(b = \frac{\pi}{6}\).} \\
\nopagebreak & & [5] & \\
5(c) & {\(V = 36\pi\left[\frac{\theta}{2} - \frac{\sin 2\theta}{4}\right]_{\pi/6}^{\pi/3}\)} & M1 & {Substitutes their limits from (b) into their integral from (a), of the form \(p\theta + q\sin 2\theta\), and subtracts.} \\
& {\(= 36\pi\left(\frac{\pi}{6} - \frac{\sqrt3}{8} - \frac{\pi}{12} + \frac{\sqrt3}{8}\right)\)} & dM1 & {Uses exact values of \(\sin\frac{2\pi}{3}\) and \(\sin\frac{\pi}{3}\), or their equality, to reach an exact expression.} \\
& {\(= 3\pi^2\)} & A1 & {cso. \(p = 3\) alone scores A0.} \\
\nopagebreak & & [3] & \\
\nopagebreak \SetCell[c=2]{r} Total for Question 5 & & 10 & \\
\end{msscheme}

\msQuestion{6}
\begin{msscheme}{
  hline{3-7} = {2-Z}{0.4pt, msGrid},
  hline{8} = {0.4pt, msGrid},
  row{8} = {font=\bfseries},
}
Question & Scheme & Marks & Notes \\
6 & {\(t^2 = y + 3 \Rightarrow x = \frac{9}{(y+3)^2}\)} & B1 & {oe. Implied by a correct integrand in \(t\).} \\
& {\(V = \pi\int_0^a \frac{81}{(y+3)^4}\,\mathrm{d}y\)} & M1 A1 & {\msNote{M1}{Uses \(\pi\int x^2\,\mathrm{d}y\) with their \(x\) in terms of \(y\), or \(\pi\int x^2\frac{\mathrm{d}y}{\mathrm{d}t}\,\mathrm{d}t\). Needs the \(\pi\); limits not needed.}\\ \msNote{A1}{Correct integral with limits 0 and \(a\), or \(162\pi\int_{\sqrt3}^{\sqrt{a+3}} t^{-7}\,\mathrm{d}t\).}} \\
& {\(\pi\left[-\frac{27}{(y+3)^3}\right]_0^a\)} & M1 A1 & {\msNote{M1}{Integrates to the form \(k(y+3)^{-3}\), or \(kt^{-6}\) in \(t\). Condone coefficient slips.}\\ \msNote{A1}{oe, such as \(-27\pi t^{-6}\).}} \\
& {\(\pi\left(1 - \frac{27}{(a+3)^3}\right) = \frac{7\pi}{8}\)} & dM1 & {Substitutes both limits into their integral, subtracts the right way round and sets equal to \(\frac{7\pi}{8}\).} \\
& {\((a+3)^3 = 216\)} & A1 & {oe, such as \(t^6 = 216\).} \\
& {\(a = 3\)} & A1 & \\
\nopagebreak \SetCell[c=2]{r} Total for Question 6 & & 8 & \\
\end{msscheme}

\msQuestion{7}
\begin{msscheme}{
  hline{3-6} = {2-Z}{0.4pt, msGrid},
  hline{7} = {0.4pt, msGrid},
  row{7} = {font=\bfseries},
}
Question & Scheme & Marks & Notes \\
7 & {\((V =) \pi\int_0^8 \left(4 - \frac{x+4}{x+1}\right)\mathrm{d}x\)} & B1 & {Needs \(\pi\) and both limits. Accept oe, such as \(32\pi - \pi\int_0^8 \frac{x+4}{x+1}\,\mathrm{d}x\).} \\
& {\(4 - \frac{x+4}{x+1} = \frac{3x}{x+1} = 3 - \frac{3}{x+1}\)\\[4pt] \(\Rightarrow \int \left(3 - \frac{3}{x+1}\right)\mathrm{d}x = \ldots\)} & M1 & {Writes their integrand as a constant plus a multiple of \(\frac{1}{x+1}\), by division or substitution, and integrates to the form \(px + q\ln(x+1)\). Condone coefficient slips.} \\
& {\(\pi\left[3x - 3\ln(x+1)\right]_0^8\)} & A1 & {Correct integral of the correct integrand, oe such as \(x + 3\ln(x+1)\) for \(\frac{x+4}{x+1}\).} \\
& {\(\pi\left((24 - 3\ln 9) - (0 - 3\ln 1)\right)\)} & dM1 & {Substitutes both limits (\(1\) and \(9\) if in \(u = x+1\)) into their integral and subtracts.} \\
& {\(= 24\pi - 6\pi\ln 3\)} & A1 & {Any exact form, such as \(3\pi(8 - \ln 9)\).} \\
\nopagebreak \SetCell[c=2]{r} Total for Question 7 & & 5 & \\
\end{msscheme}

\msQuestion{8}
\begin{msscheme}{
  hline{3-7} = {2-Z}{0.4pt, msGrid},
  hline{8} = {0.4pt, msGrid},
  row{8} = {font=\bfseries},
}
Question & Scheme & Marks & Notes \\
8 & {\(\frac{e^a - e^{-a}}{2} = 2 \Rightarrow e^{2a} - 4e^a - 1 = 0\)\\[4pt] \(\phantom{\frac{e^a}{2}} \Rightarrow e^a = \ldots\)} & M1 & {Writes \(\sinh a = 2\) in exponentials and forms a three-term quadratic in \(e^a\), or uses \(\operatorname{arsinh} x = \ln\left(x + \sqrt{x^2+1}\right)\). Condone sign slips.} \\
& {\(e^a = 2 + \sqrt5 \Rightarrow a = \ln(2 + \sqrt5)\)} & A1 & {A0 if \(2 - \sqrt5\) is retained.} \\
& {\(V = \pi\int_0^a \left(4 - \sinh^2 x\right)\mathrm{d}x\)} & B1 & {Accept the cylinder \(4\pi a\) minus \(\pi\int_0^a \sinh^2 x\,\mathrm{d}x\). Upper limit \(a\), \(\operatorname{arsinh} 2\) or their logarithm. Needs the \(\pi\).} \\
& {\(\pi\int_0^a \left(4 - \tfrac12(\cosh 2x - 1)\right)\mathrm{d}x\)\\[4pt] \(\phantom{\pi} = \pi\left[\tfrac92 x - \tfrac14\sinh 2x\right]_0^a\)} & M1 A1 & {\msNote{M1}{Uses an identity for \(\sinh^2 x\), or exponentials, and integrates to the form \(px + q\sinh 2x\) or its exponential equivalent. Condone sign and coefficient slips.}\\ \msNote{A1}{Correct integral, oe; \(\pi\) not needed.}} \\
& {\(\cosh a = \sqrt5 \Rightarrow \sinh 2a = 2(2)(\sqrt5) = 4\sqrt5\)\\[4pt] \(V = \pi\left(\tfrac92 a - \tfrac14(4\sqrt5)\right)\)} & dM1 & {Substitutes both limits and evaluates \(\sinh 2a\) exactly, from \(\sinh a = 2\) with \(\cosh^2 a = 1 + \sinh^2 a\), or from their \(e^a\).} \\
& {\(V = \pi\left(\tfrac92\ln(2 + \sqrt5) - \sqrt5\right)\)} & A1 & {Exact oe, with a logarithm; \(\operatorname{arsinh} 2\) in place of the logarithm scores A0.} \\
\nopagebreak \SetCell[c=2]{r} Total for Question 8 & & 7 & \\
\end{msscheme}

\msQuestion{9}
\begin{msscheme}{
  hline{3-6} = {2-Z}{0.4pt, msGrid},
  hline{7} = {0.4pt, msGrid},
  row{7} = {font=\bfseries},
}
Question & Scheme & Marks & Notes \\
9 & {\(V = \pi\int x^2\frac{\mathrm{d}y}{\mathrm{d}t}\,\mathrm{d}t\), \(\frac{\mathrm{d}y}{\mathrm{d}t} = 16t\)\\[4pt] \(\phantom{V} = \pi\int \left(12t(1-t)^2\right)^2(16t)\,\mathrm{d}t\)} & M1 & {Uses \(\pi\int x^2\frac{\mathrm{d}y}{\mathrm{d}t}\,\mathrm{d}t\) with their \(\frac{\mathrm{d}y}{\mathrm{d}t}\) and \(x\) substituted. Condone a slip in \(\frac{\mathrm{d}y}{\mathrm{d}t}\).} \\
& {\(V = 2304\pi\int_0^1 t^3(1-t)^4\,\mathrm{d}t\)} & A1 & {Correct integral with limits 0 and 1, oe unsimplified.} \\
& {\(2304\pi\int_0^1 \left(t^3 - 4t^4 + 6t^5 - 4t^6 + t^7\right)\mathrm{d}t\)\\[4pt] \(= 2304\pi\left[\frac{t^4}{4} - \frac{4t^5}{5} + t^6 - \frac{4t^7}{7} + \frac{t^8}{8}\right]_0^1\)} & M1 A1 & {\msNote{M1}{Expands their integrand into a polynomial, or substitutes \(u = 1 - t\), and integrates, raising at least one power by 1.}\\ \msNote{A1}{Correct integrated terms, oe.}} \\
& {\(= 2304\pi\left(\frac14 - \frac45 + 1 - \frac47 + \frac18\right)\)} & dM1 & {Substitutes \(t = 1\) and \(t = 0\), or limits in their variable, into their integrated expression.} \\
& {\(V = \frac{288\pi}{35}\ [\approx 25.9]\)} & A1 & {\(\frac{288\pi}{35}\) oe, or awrt 25.9.} \\
\nopagebreak \SetCell[c=2]{r} Total for Question 9 & & 6 & \\
\end{msscheme}

\msQuestion{10}
\begin{msscheme}{
  hline{3-8,10-11} = {2-Z}{0.4pt, msGrid},
  hline{9,12} = {0.4pt, msGrid},
  row{12} = {font=\bfseries},
}
Question & Scheme & Marks & Notes \\
10(a) & {\(V = \pi\int_{-1.088}^{1.088} \frac{8 - 6y^2 - y\sin 2y}{20}\,\mathrm{d}y\)} & B1 & {Needs \(\pi\), \(x^2\) in terms of \(y\) and limits \(\pm 1.088\); accept \(2\pi\int_0^{1.088}\) by symmetry.} \\
& {\(\int y\sin 2y\,\mathrm{d}y = -\frac{y}{2}\cos 2y\)\\[4pt] \(\phantom{\int y\sin 2y\,\mathrm{d}y =} + \int \frac12\cos 2y\,\mathrm{d}y\)} & M1 & {Uses parts on \(y\sin 2y\) to reach \(\pm py\cos 2y \pm q\int\cos 2y\,\mathrm{d}y\). Condone sign and coefficient slips.} \\
& {\(\phantom{\int y\sin 2y\,\mathrm{d}y} = -\frac{y}{2}\cos 2y + \frac14\sin 2y\)} & A1 & \\
& {\(V = \frac{\pi}{20}\left[8y - 2y^3\right.\)\\[4pt] \(\phantom{V = \frac{\pi}{20}} \left.+ \frac{y}{2}\cos 2y - \frac14\sin 2y\right]_{-1.088}^{1.088}\)} & A1 & {Correct integral of all terms, oe; the \(\frac{\pi}{20}\) may appear later.} \\
& {Substitutes \(y = 1.088\) and \(y = -1.088\)\\[4pt] and subtracts: \(V = \ldots\)} & dM1 & {Substitutes both limits into their integral of all terms and subtracts.} \\
& {\(V = 1.76\,\mathrm{cm}^3\)} & A1 & {awrt 1.76. Not for a calculator's definite integral without the integration.} \\
\nopagebreak & & [6] & \\
10(b) & {\(80 \times \text{"}1.76\text{"} = 140.8\), compared with \(140\)} & M1 & {Multiplies their (a) by 80 and compares the result with 140.} \\
& {Model uses the tallest drop, so overestimates:\\[4pt] unlikely to be enough chocolate} & E1 & {Concludes not enough, because the other drops are smaller so the total is overestimated, oe. E0 for concluding from \(140.8 > 140\) alone.} \\
\nopagebreak & & [2] & \\
\nopagebreak \SetCell[c=2]{r} Total for Question 10 & & 8 & \\
\end{msscheme}

\msQuestion{11}
\begin{msscheme}{
  hline{3-7} = {2-Z}{0.4pt, msGrid},
  hline{8} = {0.4pt, msGrid},
  row{8} = {font=\bfseries},
}
Question & Scheme & Marks & Notes \\
11 & {\(V = \pi\int x^2\frac{\mathrm{d}y}{\mathrm{d}t}\,\mathrm{d}t = \pi\int\frac{4}{(1+t)^2}\cdot 2t\,\mathrm{d}t\)} & M1 & {Forms \(\pi\int x^2\frac{\mathrm{d}y}{\mathrm{d}t}\,\mathrm{d}t\) with their \(x^2\) in \(t\) and \(\frac{\mathrm{d}y}{\mathrm{d}t} = kt\). Needs the \(\pi\); limits not needed.} \\
& {\(V = 8\pi\int\frac{t}{(1+t)^2}\,\mathrm{d}t\)} & A1 & {Correct integrand oe; limits need not be seen.} \\
& {\(\frac{t}{(1+t)^2} = \frac{1}{1+t} - \frac{1}{(1+t)^2}\)\\[4pt] \(\Rightarrow \int\ldots\,\mathrm{d}t = \ldots\)} & M1 & {Integrates by any valid method (partial fractions, \(u = 1+t\), parts) to the form \(p\ln(1+t) + \frac{q}{1+t}\) oe. Condone coefficient slips.} \\
& {\(8\pi\left[\ln(1+t) + \frac{1}{1+t}\right]\)} & A1 & {Correct integral oe, such as \(\ln(1+t) - \frac{t}{1+t}\) or in \(u\); the \(8\pi\) need not be present.} \\
& {\(t^2 + 1 = 2 \Rightarrow t = 1\), \(t^2 + 1 = 5 \Rightarrow t = 2\)\\[4pt] \(8\pi\left[\ldots\right]_1^2 = 8\pi\left(\ln 3 + \frac13 - \ln 2 - \frac12\right)\)} & dM1 & {Finds \(t = 1\) and \(t = 2\) from \(y = 2\) and \(y = 5\) and substitutes both into their integral, subtracting. Equivalent limits after a substitution accepted.} \\
& {\(V = 8\pi\ln\frac32 - \frac{4\pi}{3}\ [\approx 6.00]\)} & A1 & {oe, or awrt 6.00.} \\
\nopagebreak \SetCell[c=2]{r} Total for Question 11 & & 6 & \\
\end{msscheme}

\msQuestion{12}
\begin{msscheme}{
  hline{3-8} = {2-Z}{0.4pt, msGrid},
  hline{9} = {0.4pt, msGrid},
  row{9} = {font=\bfseries},
}
Question & Scheme & Marks & Notes \\
12 & {\(\sin x = 2\sin x\cos x \Rightarrow \cos x = \ldots\)} & M1 & {Uses \(\sin 2x = 2\sin x\cos x\) and solves for \(\cos x\), by factorising or dividing by \(\sin x\). Condone not stating \(x = 0\).} \\
& {\(x = \frac{\pi}{3}\)} & A1 & \\
& {\(V = \pi\int_0^{\pi/3}\left(\sin^2 2x - \sin^2 x\right)\mathrm{d}x\)} & B1 & {oe, such as the difference of two volume integrals. Needs the \(\pi\) and the difference of the squares.} \\
& {\(\sin^2 2x - \sin^2 x = \tfrac12(\cos 2x - \cos 4x)\)\\[4pt] \(\Rightarrow \pi\int \tfrac12(\cos 2x - \cos 4x)\,\mathrm{d}x = \ldots\)} & M1 & {Uses a double-angle identity on both squares and integrates to the form \(px + q\sin 2x + r\sin 4x\). Condone sign and coefficient slips.} \\
& {\(\pi\left[\frac{\sin 2x}{4} - \frac{\sin 4x}{8}\right]_0^{\pi/3}\)} & A1 & {oe, with or without the \(\pi\).} \\
& {\(\pi\left(\frac{\sqrt3}{8} + \frac{\sqrt3}{16}\right)\)} & dM1 & {Substitutes \(\frac{\pi}{3}\) (or their upper limit) and 0 into their integral of that form and subtracts, using exact values.} \\
& {\(V = \frac{3\sqrt3\pi}{16}\)} & A1 & {oe in terms of \(\pi\), or awrt \(0.325\pi\). A0 for awrt 1.02 alone.} \\
\nopagebreak \SetCell[c=2]{r} Total for Question 12 & & 7 & \\
\end{msscheme}

\msQuestion{13}
\begin{msscheme}{
  hline{3-5,7-12} = {2-Z}{0.4pt, msGrid},
  hline{6,13} = {0.4pt, msGrid},
  row{13} = {font=\bfseries},
}
Question & Scheme & Marks & Notes \\
13(a) & {\(\beta = \operatorname{arcosh}(1+2\alpha) \Rightarrow \cosh\beta = 1+2\alpha\)} & M1 & {Applies \(\cosh\) to the curve's equation at \(P\), or equates the log form of \(\operatorname{arcosh}(1+2\alpha)\) to \(\ln(2+\sqrt3)\).} \\
& {\(\operatorname{arcosh}2 = \ln\left(2+\sqrt{2^2-1}\right) = \ln(2+\sqrt3)\)\\[4pt] \(\Rightarrow \cosh\beta = 2\)} & M1 & {Uses \(\operatorname{arcosh}u = \ln(u+\sqrt{u^2-1})\) with \(u = 2\), or substitutes \(\beta\) into the exponential form of \(\cosh\). Not for \(\cosh\beta = 2\) from a calculator.} \\
& {\(1+2\alpha = 2 \Rightarrow \alpha = \frac12\)} & A1~AG & {Correct exact working leading without error to \(\alpha = \frac12\). Working back from \(\alpha = \frac12\) earns nothing.} \\
\nopagebreak & & [3] & \\
13(b) & {\(\cosh y = 1+2x \Rightarrow x = \frac{\cosh y - 1}{2}\)} & B1 & \\
& {\(V = \pi\int_0^{\beta}\left(\frac{\cosh y - 1}{2}\right)^2\mathrm{d}y\)} & B1 & {Needs \(\pi\), \(x^2\) as their function of \(y\), and limits \(0\) and \(\beta\) oe.} \\
& {\((\cosh y - 1)^2 = \tfrac12(\cosh 2y + 1) - 2\cosh y + 1\)\\[4pt] \(\Rightarrow \int (\cosh y - 1)^2\,\mathrm{d}y = \ldots\)} & M1 & {Expands, uses a double-angle identity or exponentials, and integrates to the form \(p\sinh 2y + q\sinh y + ry\) or its exponential equivalent. Condone sign and coefficient slips.} \\
& {\(V = \frac{\pi}{4}\left[\tfrac14\sinh 2y - 2\sinh y + \tfrac32 y\right]_0^{\beta}\)} & A1 & {Correct integral oe, including exponential forms. The factor \(\frac{\pi}{4}\) may be absent here.} \\
& {\(\cosh\beta = 2 \Rightarrow \sinh\beta = \sqrt3,\ \sinh 2\beta = 4\sqrt3\)\\[4pt] \(V = \frac{\pi}{4}\left(\tfrac14(4\sqrt3) - 2\sqrt3 + \tfrac32\beta\right)\)} & dM1 & {Substitutes the limits, with \(\sinh\beta\) and \(\sinh 2\beta\) (or \(e^{\pm\beta}\), \(e^{\pm 2\beta}\)) found exactly. Condone slips. Not for decimals.} \\
& {\(V = \frac{\pi}{8}\left(3\ln(2+\sqrt3) - 2\sqrt3\right)\)} & A1 & {Any exact equivalent; \(\ln(2+\sqrt3)\) may be left as \(\beta\) or \(\operatorname{arcosh}2\).} \\
\nopagebreak & & [6] & \\
\nopagebreak \SetCell[c=2]{r} Total for Question 13 & & 9 & \\
\end{msscheme}

\msQuestion{14}
\begin{msscheme}{
  hline{3-6,8-13} = {2-Z}{0.4pt, msGrid},
  hline{7,14} = {0.4pt, msGrid},
  row{14} = {font=\bfseries},
}
Question & Scheme & Marks & Notes \\
14(a) & {\(x = \sin\theta - \frac13\left(3\sin\theta - 4\sin^3\theta\right)\)} & M1 & {Uses a triple-angle identity for \(\sin 3\theta\), or any valid expansion, to write \(x\) in terms of \(\sin\theta\) alone. Condone sign slips.} \\
& {\(x = \frac43\sin^3\theta\)} & A1 & \\
& {\(\cos\theta = 1 - y \Rightarrow \sin^2\theta = 1 - (1-y)^2\)\\[4pt] \(x^2 = \text{"}\frac{16}{9}\text{"}\sin^6\theta\)\\[4pt] \(\phantom{x^2} = \text{"}\frac{16}{9}\text{"}\left(1 - (1-y)^2\right)^3\)} & M1 & {Uses \(\sin^2\theta + \cos^2\theta = 1\) with \(\cos\theta = 1 - y\), and substitutes into their \(x^2\) of the form \(k\sin^6\theta\).} \\
& {\(x^2 = \frac{16}{9}\left(2y - y^2\right)^3 = \frac{16}{9}y^3(2-y)^3\)} & A1~AG & {Correct substituted equation leading without error to the printed result; combined steps accepted.} \\
\nopagebreak & & [4] & \\
14(b) & {\(V = \pi\int_0^2 \frac{16}{9}y^3(2-y)^3\,\mathrm{d}y\)} & B1 & {Needs the \(\pi\), the integrand from (a) and limits 0 and 2.} \\
& {\(\int_0^2 \left(8y^3 - 12y^4 + 6y^5 - y^6\right)\mathrm{d}y\)\\[4pt] \(\phantom{\int} = \left[2y^4 - \frac{12}{5}y^5 + y^6 - \frac17 y^7\right]_0^2\)} & M1 A1 & {\msNote{M1}{Integrates \(y^3(2-y)^3\) by expanding or any valid method, raising the power of at least one term by 1. Condone expansion slips.}\\ \msNote{A1}{Correct integral, with or without the factor \(\frac{16\pi}{9}\).}} \\
& {\(\frac{16\pi}{9}\left(32 - \frac{384}{5} + 64 - \frac{128}{7}\right)\)} & dM1 & {Substitutes \(y = 2\) and \(y = 0\) into their integrated expression and subtracts.} \\
& {\(\frac{512\pi}{315}\ [\approx 5.11]\)} & A1 & {oe, or awrt 5.11.} \\
& {\(\left(\frac32\right)^3 \times \text{"}\frac{512\pi}{315}\text{"}\)} & M1 & {Multiplies their model volume by the cube of 3 over their model height, or integrates with lengths scaled by \(\frac32\).} \\
& {\(\frac{192\pi}{35} \approx 17.2\,\mathrm{cm}^3\)} & A1 & {oe, or awrt 17.2.} \\
\nopagebreak & & [7] & \\
\nopagebreak \SetCell[c=2]{r} Total for Question 14 & & 11 & \\
\end{msscheme}

\msQuestion{15}
\begin{msscheme}{
  hline{3-4,6-11} = {2-Z}{0.4pt, msGrid},
  hline{5,12} = {0.4pt, msGrid},
  row{12} = {font=\bfseries},
}
Question & Scheme & Marks & Notes \\
15(a) & {\(\cos 5\theta = \cos(4\theta + \theta) = \cos 4\theta\cos\theta - \sin 4\theta\sin\theta\)\\[4pt] \(\cos 3\theta = \cos(4\theta - \theta) = \cos 4\theta\cos\theta + \sin 4\theta\sin\theta\)} & M1 A1 & {\msNote{M1}{Expands \(\cos(4\theta + \theta)\) and \(\cos(4\theta - \theta)\) by compound-angle formulae. Condone sign slips.}\\ \msNote{A1}{Both expansions correct.}} \\
& {\(\cos 5\theta + \cos 3\theta = 2\cos 4\theta\cos\theta\)\\[4pt] \(\Rightarrow \cos 4\theta\cos\theta \equiv \tfrac12(\cos 5\theta + \cos 3\theta)\)} & A1~AG & {Correct expansions combined without error to the printed identity, with LHS = RHS or a conclusion in words.} \\
\nopagebreak & & [3] & \\
15(b) & {\(V = \pi\int x^2\frac{\mathrm{d}y}{\mathrm{d}\theta}\,\mathrm{d}\theta\), \(\frac{\mathrm{d}y}{\mathrm{d}\theta} = 24\cos\theta\)} & M1 & {Uses \(\pi\int x^2\,\mathrm{d}y\) with \(x^2\) in terms of \(\theta\) and \(\mathrm{d}y\) as their \(k\cos\theta\,\mathrm{d}\theta\). Limits not needed.} \\
& {\(V = 5400\pi\int_0^{\pi/2}\sin^2 2\theta\cos\theta\,\mathrm{d}\theta\)} & A1 & {Correct integral with limits \(0\) and \(\frac{\pi}{2}\), oe, need not be simplified.} \\
& {\(2700\pi\int_0^{\pi/2}\left(\cos\theta - \tfrac12\cos 5\theta - \tfrac12\cos 3\theta\right)\mathrm{d}\theta\)} & M1 & {Simplifies the integrand, by \(\sin^2 2\theta = \frac12(1 \pm \cos 4\theta)\) and part (a) or otherwise, and integrates at least one term correctly. Condone sign slips.} \\
& {\(= 2700\pi\left[\sin\theta - \tfrac{1}{10}\sin 5\theta - \tfrac16\sin 3\theta\right]_0^{\pi/2}\)} & A1 & {Correct integral, oe, including \(5400\pi\left[\frac43u^3 - \frac45u^5\right]_0^1\) with \(u = \sin\theta\).} \\
& {\(= 2700\pi\left[\left(1 - \tfrac{1}{10} + \tfrac16\right) - 0\right]\)} & dM1 & {Substitutes both limits, or the matching limits in a new variable, into their integral and subtracts.} \\
& {\(V = 2880\pi\ (\mathrm{mm}^3)\)} & A1 & {Or awrt 9050.} \\
\nopagebreak & & [6] & \\
\nopagebreak \SetCell[c=2]{r} Total for Question 15 & & 9 & \\
\end{msscheme}

\end{document}
